\clearpage\selectlanguage{english}\fancyhead[L]{\small NACHLASS.}\pdfbookmark[0]{DETERMINATIO SERIEI NOSTRAE PER AEQUATIONEM DIFFERENTIALEM SECUNDI ORDINIS.}{nachlass216}
\SourceLeaf{216}
\begin{center}
\textbf{N A C H L A S S.}\par\medskip
\textbf{DETERMINATIO SERIEI NOSTRAE}\par\medskip
\textbf{PER AEQUATIONEM DIFFERENTIALEM SECUNDI ORDINIS.}\par\medskip
\rule{0.08\linewidth}{0.3pt}\par\bigskip
\end{center}
\SixthArticle{38}
Statuendo brevitatis causa \(F(\alpha,\beta,\gamma,x)=P\), habemus per art. 4
\[
\frac{\mathrm{d}P}{\mathrm{d}x}=\frac{\alpha\beta}{\gamma}
F(\alpha+1,\beta+1,\gamma+1,x)
\]
atque hinc differentiando denuo
\[
\frac{\mathrm{dd}P}{\mathrm{d}x^2}
=\frac{\alpha\beta(\alpha+1)(\beta+1)}{\gamma(\gamma+1)}
F(\alpha+2,\beta+2,\gamma+2,x)
\]
Hinc aequatio IX art. 10 suppeditat
\[
\text{[80]}\qquad
0=\alpha\beta P-(\gamma-(\alpha+\beta+1)x)
\frac{\mathrm{d}P}{\mathrm{d}x}-(x-xx)\frac{\mathrm{dd}P}{\mathrm{d}x^2}
\]
Haecce aequatio differentio-differentialis tamquam definitio exactior
functionis nostrae considerari potest; sed quoniam \(P=F(\alpha,\beta,\gamma,x)\)
non est integrale completum sed particulare tantum (quod constantes non
accesserunt), adiicere oportet conditionem, ut \(P\) incipiat a valore 1
pro \(x=0\) simulque pro eodem valore ipsius \(x\) supponatur
\(\frac{\mathrm{d}P}{\mathrm{d}x}=\frac{\alpha\beta}{\gamma}\)
atque
\(\frac{\mathrm{dd}P}{\mathrm{d}x^2}
=\frac{\alpha\beta(\alpha+1)(\beta+1)}{\gamma(\gamma+1)}\).

Ita pro quovis valore ipsius \(x\), ad quem a valore \(x=0\) per gradus
continuos transiisti, ita tamen, ut valorem \(x=1\), pro quo \(x-xx=0\),
non attigeris, \(P\) erit quantitas perfecte determinata; sed manifesto
\emph{hoc modo} ad valores reales ipsius \(x\) positivos unitate majores
pervenire nequis nisi transeundo per va\SourceLeaf{217}%
lores imaginarios, quod quum infinitis modis diversis absque continuitatis
praejudicio fieri possit, hinc nondum liquet, annon eidem valori ipsius
\(x\) plures, quin adeo infinite multi valores discreti ipsius \(P\)
respondeant, sicuti in pluribus functionibus transscendentibus magis notis
evenire constat. Sed de hoc argumento in posterum fusius loqui nobis
reservamus, quum hoc loco casus is potissimum, ubi \(x\) accipitur infra
vel saltem non ultra unitatem positivam, atque \(P\) aequalis summae
seriei \(F(\alpha,\beta,\gamma,x)\) tractetur.

\SixthArticle{39}
Scribendo in aequatione 80, \(1-y\) pro \(x\), transit ea in hanc
\[
\begin{aligned}&0=\alpha\beta P-(\alpha+\beta+1-\gamma-(\alpha+\beta+1)y)
\frac{\mathrm{d}P}{\mathrm{d}y}\\&-(y-yy)\frac{\mathrm{dd}P}{\mathrm{d}y^2}\end{aligned}
\]
quae habet formam similem ut illa. Hinc statim prodit aliud integrale particulare
\[
\begin{aligned}&P=F(\alpha,\beta,\alpha+\beta+1-\gamma,y)\\&=F(\alpha,\beta,\alpha+\beta+1-\gamma,1-x)\end{aligned}
\]
unde per principia nota sequitur integrale completum aequationis 80,
\[
\begin{aligned}&\text{[81]}\qquad
P=MF(\alpha,\beta,\gamma,x)\\&+NF(\alpha,\beta,\alpha+\beta+1-\gamma,1-x)\end{aligned}
\]
denotantibus \(M,N\) constantes arbitrarias.

Ceterum obiter hic observamus, ad formam aequationis 80 facile reduci
posse aequationem generaliorem
\[
0=AP+(B+Cy)\frac{\mathrm{d}P}{\mathrm{d}y}
+(D+Ey+Fyy)\frac{\mathrm{dd}P}{\mathrm{d}y^2}
\]
Sunto enim radices aequationis \(0=D+Ey+Fyy\) hae.\ \(y=a,y=b\),
sive \(D+Ey+Fyy\) indefinite aequalis producto \(F(y-a)(y-b)\)
patetque statuendo \(\frac{y-a}{b-a}=x\) atque determinando
\(\alpha,\beta,\gamma\) ita ut fiat
\[
\alpha\beta=\frac{A}{F},\quad
\alpha+\beta+1=\frac{C}{F},\quad
\gamma=-\frac{B+aC}{F(b-a)}
\]
illam transire in aequationem 80.

\SixthArticle{40}
Adiumento aequationis differentialis 80 eruere licet complura theoremata
maxime memorabilia circa seriem nostram, tum generalia tum magis specialia,
neque dubitamus, quin multa plura atque graviora adhuc lateant, ulterioribus
curis servata. Quae hactenus nobis revelare contigit, hic in conspectum
producamus.
\par\SourceLeaf{218}
Statuamus $P=(1-x)^\mu P'$, eritque
\[
\begin{aligned}
\frac{\mathrm{d}P}{\mathrm{d}x}&=-\mu(1-x)^{\mu-1}P'+(1-x)^\mu\frac{\mathrm{d}P'}{\mathrm{d}x}\\
\frac{\mathrm{dd}P}{\mathrm{d}x^2}&=\mu(\mu-1)(1-x)^{\mu-2}P'-2\mu(1-x)^{\mu-1}\frac{\mathrm{d}P'}{\mathrm{d}x}\\&\qquad+(1-x)^\mu\frac{\mathrm{dd}P'}{\mathrm{d}x^2}
\end{aligned}
\]
Quibus valoribus in aequatione 80 substitutis prodit dividendo per $(1-x)^{\mu-1}$
\[
\begin{aligned}
0={}&P'\{\alpha\beta(1-x)+(\gamma-(\alpha+\beta+1)x)\mu-x(\mu\mu-\mu)\}\\
&-\frac{\mathrm{d}P'}{\mathrm{d}x}\{(\gamma-(\alpha+\beta+1)x)-2\mu x\}(1-x)\\&-\frac{\mathrm{dd}P'}{\mathrm{d}x^2}\{x-xx\}(1-x)
\end{aligned}
\]
Determinemus $\mu$ ita, ut multiplicator ipsius $P'$ per $1-x$ divisibilis evadat, quod fiet vel statuendo $\mu=0$ vel $\mu=\gamma-\alpha-\beta$. Suppositio prior nihil novi doceret, sed valor posterior substitutus producit
\[
\begin{aligned}&0=P'\{\alpha\beta-\alpha\gamma-\beta\gamma+\gamma\gamma\}-\frac{\mathrm{d}P'}{\mathrm{d}x}\{\gamma-(2\gamma-\alpha-\beta+1)x\}\\&-\frac{\mathrm{dd}P'}{\mathrm{d}x^2}\{x-xx\}\end{aligned}
\]
sive
\[
\begin{aligned}&0=P'(\gamma-\alpha)(\gamma-\beta)-\frac{\mathrm{d}P'}{\mathrm{d}x}\{\gamma-((\gamma-\alpha)+(\gamma-\beta)+1)x\}\\&-\frac{\mathrm{dd}P'}{\mathrm{d}x^2}(x-xx)\end{aligned}
\]
quae prorsus eandem formam habet ut aequatio 80. Quare quum pro $x=0$, manifesto fiat $P'=1$ atque $\frac{\mathrm{d}P'}{\mathrm{d}x}=\frac{\alpha\beta}{\gamma}-\mu=\frac{(\gamma-\alpha)(\gamma-\beta)}{\gamma}$, patet ipsius integrale esse $P'=F(\gamma-\alpha,\gamma-\beta,\gamma,x)$, ita ut generaliter habeatur
\[
\text{[82]}\qquad F(\gamma-\alpha,\gamma-\beta,\gamma,x)=(1-x)^{\alpha+\beta-\gamma}F(\alpha,\beta,\gamma,x)
\]
Hinc petenda est transformatio seriei
\[
1+\frac{2\,.\,8}{9}x+\frac{3\,.\,8\,.\,10}{9\,.\,11}xx+\frac{4\,.\,8\,.\,10\,.\,12}{9\,.\,11\,.\,13}x^3+\text{etc.}=F(2,4,\tfrac{9}{2},x)
\]
in
\[
\begin{aligned}&(1-x)^{-\frac{3}{2}}\left(1+\frac{1\,.\,5}{2\,.\,9}x+\frac{1\,.\,3\,.\,5\,.\,7}{2\,.\,4\,.\,9\,.\,11}xx+\text{etc.}\right)\\&=(1-x)^{-\frac{3}{2}}F(\tfrac{5}{2},\tfrac{1}{2},\tfrac{9}{2},x)\end{aligned}
\]
quam in Ephemeridibus Astronomicis Berolinensibus 1814 p. 257 [Zusatz zu Art. 90 und 100 der Theoria motus] sine demonstratione indicaveramus.

\SixthArticle{41}

Statuamus porro $P=x^\mu P'$, ita ut fiat
\[
\begin{aligned}
\frac{\mathrm{d}P}{\mathrm{d}x}&=\mu x^{\mu-1}P'+x^\mu\frac{\mathrm{d}P'}{\mathrm{d}x}\\
\frac{\mathrm{dd}P}{\mathrm{d}x^2}&=(\mu\mu-\mu)x^{\mu-2}P'+2\mu x^{\mu-1}\frac{\mathrm{d}P'}{\mathrm{d}x}+x^\mu\frac{\mathrm{dd}P'}{\mathrm{d}x^2}
\end{aligned}
\]
\SourceLeaf{219}
quibus valoribus in 80 substitutis, fit dividendo per $x^{\mu-1}$
\[
\begin{aligned}
0={}&P'\{\alpha\beta x-(\gamma-(\alpha+\beta+1)x)\mu-(1-x)(\mu\mu-\mu)\}\\
&-\frac{\mathrm{d}P'}{\mathrm{d}x}\{\gamma-(\alpha+\beta+1)x+2\mu(1-x)\}x\\
&-\frac{\mathrm{dd}P'}{\mathrm{d}x^2}(xx-x^3)
\end{aligned}
\]
Multiplicator ipsius $P'$ in hac formula fit divisibilis per $x$ statuendo $\mu=0$ vel $\mu=1-\gamma$; valor posterior producit
\[
\begin{aligned}
0={}&P'(\alpha\beta+\alpha+\beta+1-2\gamma-\alpha\gamma-\beta\gamma+\gamma\gamma)\\
&-\frac{\mathrm{d}P'}{\mathrm{d}x}(2-\gamma-(\alpha+\beta+3-2\gamma)x)\\
&-\frac{\mathrm{dd}P'}{\mathrm{d}x^2}(x-xx)
\end{aligned}
\]
Comparando hanc aequationem cum 80, cuius forma prorsus similis est, patet quae illic fuerant $P$, $\alpha$, $\beta$, $\gamma$ hic esse $P'$, $\alpha+1-\gamma$, $\beta+1-\gamma$, $2-\gamma$: quare quum illius integrale completum assignaverimus, manifestum est, $P'$ contentam fore sub formula
\[
\begin{aligned}
P'={}&MF(\alpha+1-\gamma,\beta+1-\gamma,2-\gamma,x)\\
&+NF(\alpha+1-\gamma,\beta+1-\gamma,\alpha+\beta+1-\gamma,1-x)
\end{aligned}
\]
denotantibus $M$, $N$ quantitates constantes, sive
\[
\begin{aligned}
&\text{[83]}\qquad F(\alpha,\beta,\gamma,x)\\&=Mx^{1-\gamma}F(\alpha+1-\gamma,\beta+1-\gamma,2-\gamma,x)\\
&+Nx^{1-\gamma}F(\alpha+1-\gamma,\beta+1-\gamma,\alpha+\beta+1-\gamma,1-x)
\end{aligned}
\]
ubi constantes $M$, $N$ ab elementis $\alpha$, $\beta$, $\gamma$ pendebunt.

\SixthArticle{42}

Ex aequatione 82 sequitur
\[
\begin{aligned}
&F(\alpha+1-\gamma,\beta+1-\gamma,2-\gamma,x)\\&=(1-x)^{\gamma-\alpha-\beta}F(1-\alpha,1-\beta,2-\gamma,x)\\[5pt]
&F(\alpha+1-\gamma,\beta+1-\gamma,\alpha+\beta+1-\gamma,1-x)\\&=x^{\gamma-1}F(\alpha,\beta,\alpha+\beta+1-\gamma,1-x)
\end{aligned}
\]
unde statuendo
\[
\frac{1}{N}=f(\alpha,\beta,\gamma),\qquad-\frac{M}{N}=g(\alpha,\beta,\gamma)
\]
\SourceLeaf{220}
aequatio 83 fit
\[
\begin{aligned}
&F(\alpha,\beta,\alpha+\beta+1-\gamma,1-x)\\&=f(\alpha,\beta,\gamma)F(\alpha,\beta,\gamma,x)\\
&\quad+g(\alpha,\beta,\gamma)(1-x)^{\gamma-\alpha-\beta}x^{1-\gamma}F(1-\alpha,1-\beta,2-\gamma,x)
\end{aligned}
\]
Eidem aequationi adiumento formulae 82 hanc quoque formam tribuere licet
\[
\begin{aligned}
&x^{1-\gamma}F(\alpha+1-\gamma,\beta+1-\gamma,\alpha+\beta+1-\gamma,1-x)=\\
&f(\alpha,\beta.\gamma)(1-x)^{\gamma-\alpha-\beta}F(\gamma-\alpha,\gamma-\beta,\gamma,x)
\\&+g(\alpha,\beta,\gamma)x^{1-\gamma}F(\alpha+1-\gamma,\beta+1-\gamma,2-\gamma,x)
\end{aligned}
\]
sive dividendo per \(x^{1-\gamma}\), mutandoque resp.\ \(\alpha,\beta,\gamma\) in \(\alpha+1-\gamma,\beta+1-\gamma,2-\gamma\)
\[
\begin{aligned}
&F(\alpha,\beta,\alpha+\beta+1-\gamma,1-x)\\
&\quad=g(\alpha+1-\gamma,\beta+1-\gamma,2-\gamma)F(\alpha,\beta,\gamma,x)\\
&\qquad+f(\alpha+1-\gamma,\beta+1-\gamma,2-\gamma)\\&\qquad(1-x)^{\gamma-\alpha-\beta}x^{1-\gamma}F(1-\alpha,1-\beta,2-\gamma,x)
\end{aligned}
\]
Quae quum identica esse debeat cum formula praecedenti, habemus
\[
g(\alpha,\beta,\gamma)=f(\alpha+1-\gamma,\beta+1-\gamma,2-\gamma)
\]
itaque
\[
\begin{aligned}
\text{[84]}\quad &F(\alpha,\beta,\alpha+\beta+1-\gamma,1-x)\\
&=f(\alpha,\beta,\gamma)F(\alpha,\beta,\gamma,x)\\
&\quad+f(\alpha+1-\gamma,\beta+1-\gamma,2-\gamma)\\&\qquad(1-x)^{\gamma-\alpha-\beta}x^{1-\gamma}F(1-\alpha,1-\beta,2-\gamma,x)
\end{aligned}
\]
\SixthArticle{43}
Iam ut indolem functionis \(f(\alpha,\beta,\gamma)\) eruamus, statuamus \(x=0\). Tunc patet, esse \(F(\alpha,\beta,\gamma,x)=1.\ x^{1-\gamma}=0\), quoties quidem \(1-\gamma\) fuerit quantitas positiva. Sed per aequationem 48 habemus
\[
F(\alpha,\beta,\alpha+\beta+1-\gamma,1)=\frac{\Pi(\alpha+\beta-\gamma)\Pi(-\gamma)}{\Pi(\alpha-\gamma)\Pi(\beta-\gamma)}
\]
Quare sub eadem restrictione demonstratum est, fieri
\[
\text{[85]}\qquad f(\alpha,\beta,\gamma)=\frac{\Pi(\alpha+\beta-\gamma)\Pi(-\gamma)}{\Pi(\alpha-\gamma)\Pi(\beta-\gamma)}
\]
Hanc vero formulam generalem esse, ita demonstramus.
\par\SourceLeaf{221}
Differentiando aequationem 84 provenit
\[
\begin{aligned}
&-\frac{\alpha\beta}{\alpha+\beta+1-\gamma}F(\alpha+1,\beta+1,\alpha+\beta+2-\gamma,1-x)\\
&=\frac{\alpha\beta}{\gamma}f(\alpha,\beta,\gamma)F(\alpha+1,\beta+1,\gamma+1.\ x)\\
&\quad+f(\alpha+1-\gamma,\beta+1-\gamma,2-\gamma)(1-x)^{\gamma-\alpha-\beta-1}x^{-\gamma}\\&\qquad\Bigl\{((1-\gamma)(1-x)-(\gamma-\alpha-\beta)x)F(1-\alpha,1-\beta,2-\gamma,x)\\
&\qquad+\frac{(1-\alpha)(1-\beta)}{2-\gamma}(x-xx)F(2-\alpha,2-\beta,3-\gamma,x)\Bigr\}
\end{aligned}
\]
Sed per formulam IX art.\ 10 fit mutando \(\alpha,\beta,\gamma\) in \(-\alpha,-\beta,1-\gamma\)
\[
\begin{aligned}
&(1-\gamma)(2-\gamma)F(-\alpha,-\beta,1-\gamma,x)\\&=(2-\gamma)(1-\gamma+(\alpha+\beta-1)x)F(1-\alpha,1-\beta,2-\gamma,x)\\
&\quad+(1-\alpha)(1-\beta)(x-xx)F(2-\alpha,2-\beta,3-\gamma,x)
\end{aligned}
\]
unde aequatio praecedens transit in hanc
\[
\begin{aligned}
&F(\alpha+1,\beta+1,\alpha+\beta+2-\gamma,1-x)\\
&=-\frac{\alpha+\beta+1-\gamma}{\gamma}f(\alpha,\beta,\gamma)F(\alpha+1,\beta+1,\gamma+1.\ x)\\
&\quad-\frac{(\alpha+\beta+1-\gamma)(1-\gamma)}{\alpha\beta}f(\alpha+1-\gamma,\beta+1-\gamma,2-\gamma)\\&\qquad(1-x)^{\gamma-\alpha-\beta-1}x^{-\gamma}F(-\alpha,-\beta,1-\gamma,x)
\end{aligned}
\]
Mutando autem in aequatione 84, \(\alpha,\beta,\gamma\) in \(\alpha+1,\beta+1,\gamma+1\) fit
\[
\begin{aligned}
&F(\alpha+1,\beta+1,\alpha+\beta+2-\gamma,1-x)\\
&\quad=f(\alpha+1,\beta+1,\gamma+1)F(\alpha+1,\beta+1,\gamma+1,x)\\
&\qquad+f(\alpha+1-\gamma,\beta+1-\gamma,1-\gamma)\\&\qquad(1-x)^{\gamma-\alpha-\beta-1}x^{-\gamma}F(-\alpha,-\beta,1-\gamma,x)
\end{aligned}
\]
Quare quum facile perspiciatur, has duas aequationes identicas esse debere, fit generaliter
\[
f(\alpha+1,\beta+1,\gamma+1)=\frac{\alpha+\beta+1-\gamma}{-\gamma}f(\alpha,\beta,\gamma)
\]
sive mutando \(\alpha,\beta,\gamma\) in \(\alpha-1,\beta-1,\gamma-1\)
\[
\begin{aligned}
f(\alpha,\beta,\gamma)&=\frac{\alpha+\beta-\gamma}{1-\gamma}f(\alpha-1,\beta-1,\gamma-1)\\
&=\frac{\alpha+\beta-\gamma\,.\,\alpha+\beta-\gamma-1}{1-\gamma\,.\,2-\gamma}f(\alpha-2,\beta-2,\gamma-2)
\end{aligned}
\]
etc.\ unde facile concluditur, esse generaliter pro quovis valore integro ipsius \(k\)
\[
f(\alpha,\beta,\gamma)=\frac{\Pi(\alpha+\beta-\gamma)\Pi(-\gamma)}{\Pi(\alpha+\beta-\gamma-k)\Pi(k-\gamma)}.f(\alpha-k,\beta-k,\gamma-k)
\]
